\section{引言}
\label{sec:introduction}

连续情感回归用于估计观众观看视频时逐时变化的效价与唤醒度。准确的群体层面情感轨迹可以支持视频理解，并为视频剪辑、广告投放、内容推荐或新观众预期反应的粗略初始化提供基线。本文聚焦这些用途的预测前提：降低新观众观看熟悉、时间对齐视频时的平均绝对误差（MAE）。

EEG 提供毫秒级电生理测量，功能性近红外光谱（fNIRS）则捕获较慢的血流动力学变化。二者的互补特性推动了 EEG--fNIRS 情感模型的发展
\cite{bandara2018fnirs,sun2020multimodal,chen2024eegfnirs}。然而，跨参与者预测仍然困难，因为解剖结构、传感器接触、基线状态和反应方式可能产生比目标情感变化更大的偏移
\cite{li2022eegreview,li2018crosssubject}。与使用视频已有信息相比，为每位新观众采集两种模态还会带来显著更高的传感和处理成本。

熟悉视频场景改变了可用的信息。自然电影能够诱发 BOLD 与 EEG 反应中时间锁定的跨参与者相似性
\cite{hasson2004intersubject,kauppi2010isc,dmochowski2012correlated}。
已有研究还考察了这些神经反应的可靠性与情绪、注意和观众偏好之间的关系
\cite{nummenmaa2012synchrony,ki2016attention,dmochowski2014audience}。
当训练参与者和新观众观看相同且时间对齐的视频时，匹配的标签可以形成由视频身份和视频内时间索引的群体情感轨迹。有效预测器必须在每个训练划分内重建该轨迹，防止留出标签进入估计过程。

本文将共享反应操作化为折内视频—时间先验，将其与 EEG--fNIRS 分支配对，并通过偏重先验的固定融合加以组合。研究围绕三个问题展开：

\begin{description}
  \item[RQ1.] 视频—时间先验与固定融合在熟悉视频新观众预测中达到怎样的 MAE？
  \item[RQ2.] 全局水平、视频身份、视频内时间或 EEG--fNIRS 中的哪一种信息来源解释了 MAE 的降低？
  \item[RQ3.] 剩余增益如何随参与者、视频和视频内时间而变化？
\end{description}

本文的贡献有三点：

\begin{itemize}
  \item 我们形式化了面向熟悉视频新观众、以 MAE 为目标的连续情感预测任务，并构建一种折内视频—时间先验；该先验排除被评估参与者的标签，且边际推理成本较低。
  \item 我们通过实证表明，视频—时间先验本身就是一个强预测器，而偏重先验的固定融合与 EEG--fNIRS 分支组合后，在内部参与者留出和外部参与者独立评估中均取得最低 MAE。
  \item 我们通过来源显式消融和成对残差分析，分离全局水平、视频身份、视频内时间和 EEG--fNIRS 修正，揭示生理贡献较小且具有异质性。
\end{itemize}

\section{相关工作}
\label{sec:related}

\subsection{生理情感回归}

生理情感识别通常在离散类别或效价—唤醒度空间中表征情感
\cite{russell1980circumplex,koelstra2012deap}。EEG 研究已识别出有信息量的频谱与空间模式
\cite{zheng2015critical}，而 fNIRS 研究使用氧合相关统计量建模情感状态
\cite{bandara2018fnirs}。多模态系统通过图卷积、注意力或表征对齐组合两类信号
\cite{sun2020multimodal,chen2024eegfnirs,fang2026asac}。这些机制构成本文 EEG--fNIRS 分支的设计依据。

跨参与者差异仍是一个长期存在的障碍
\cite{li2022eegreview,li2018crosssubject}。静息记录和迁移学习可以对齐参与者或个性化情感解码器
\cite{morioka2015resting,zheng2016personalizing}。这些方法处理参与者偏移，却没有量化熟悉视频的身份和播放时间中已经包含的低成本信息。因此，对该场景进行公平分析，需要在生理模型旁设置显式的视频—时间基线。

\subsection{自然刺激与共享反应}

自然电影能够诱发 BOLD 与 EEG 反应中时间锁定的跨参与者相似性
\cite{hasson2004intersubject,kauppi2010isc,dmochowski2012correlated}。
已有研究从情绪、注意和观众偏好的角度考察这些神经反应的可靠性
\cite{nummenmaa2012synchrony,ki2016attention,dmochowski2014audience}。
这些工作建立了共享的刺激锁定神经结构，但并未直接建立连续情感标签预测器。

在本文的重复刺激场景中，训练参与者的匹配标签支持另一种操作化方式：由视频身份和视频内时间索引的折内群体轨迹。视频—时间先验将该轨迹转化为低成本回归基线，并与 EEG--fNIRS 分支直接比较。

\subsection{评估、先验与泄漏控制}

评估协议决定了报告分数能够支持哪一种泛化结论。分组交叉验证可防止参与者直接重叠，但各折重复使用视频时，评估的仍是新参与者在熟悉内容上的表现
\cite{saeb2017usecase,roberts2017crossvalidation}。视频留出设计和参与者—视频交叉设计回答不同问题。在评估折内进行模型选择还可能使估计产生偏差
\cite{varma2006bias,varoquaux2017decoders}。

对于熟悉视频预测，只有当视频—时间先验由训练参与者构建，并应用于完全相同的留出样本键时，它才是有效的。为此，我们在每个内部折内重建先验，并使所有外部目标始终处于训练与先验构建之外。随后，来源显式变体分离全局情感水平、视频身份、视频内时间和 EEG--fNIRS 测量所提供的信息。

\section{方法}
\label{sec:method}

\subsection{问题定义与流程}

对于参与者 \(s\)、熟悉视频 \(v\) 和秒级时刻 \(t\)，目标
\(y_{s,v,t}\in[1,255]^2\) 包含效价与唤醒度。首要评价指标是两个情感维度上的样本汇总 MAE：
\begin{equation}
 \operatorname{MAE} =
 \frac{1}{2N}\sum_{i=1}^{N}\sum_{d\in\{V,A\}}
 \left|y_i^d-\hat y_i^d\right|.
\end{equation}

该流程在融合前产生两个互补估计。视频—时间先验 \(\bar p_{v,t}\) 检索与熟悉视频及秒级时刻对应的群体轨迹。EEG--fNIRS 分支 \(m_{s,v,t}\) 提取一秒通道特征，使用图网络分别编码两种模态，通过注意力交换跨模态信息，并从池化 token 回归效价与唤醒度。固定凸组合输出最终预测 \(\hat y_{s,v,t}\)。该设计使低成本共享反应与参与者特异性生理修正能够被分别测量。

\subsection{折内视频—时间先验}

在每个训练折内，我们对每个视频—时间坐标的标签计算跨参与者中位数：
\begin{equation}
 p^d_{v,t} =
 \operatorname{median}_{s\in\mathcal S_{\mathrm{tr}}} y^d_{s,v,t},
 \qquad d\in\{V,A\}.
\end{equation}
随后使用半径为三的移动平均对每条轨迹进行平滑：
\begin{equation}
 \bar p^d_{v,t} =
 \frac{1}{|\mathcal I_{v,t}|}
 \sum_{u\in\mathcal I_{v,t}} p^d_{v,u},
 \quad
 \mathcal I_{v,t}=\{u:|u-t|\leq3,\ 0\leq u<T_v\}.
\end{equation}
平滑从不跨越视频边界。中位数限制个体轨迹的影响，折内重建则防止留出标签决定其自身预测。推理时，先验仅按视频身份和秒级时刻查表，因此训练轨迹构建完成后，其边际输入与计算成本较低。

\subsection{EEG--fNIRS 分支}

EEG--fNIRS 分支首先构建通道级特征，随后依次进行模态内图编码、双向跨模态注意力和池化回归。信号使用相应试验的基线进行中心化，EEG 被重采样至 200 Hz。对于每个 EEG 通道的一秒片段，我们计算 1--45 Hz 六个频带中的相对对数功率和微分熵，以及 Hjorth 活动度、移动性和复杂度
\cite{hjorth1970eeg}。拼接 \(t-1\)、\(t\) 与 \(t+1\) 的特征后，得到 \(64\times45\) 的 EEG 张量。

对于每个 fNIRS 通道和六类信号，我们计算均值、标准差、斜率、偏度和峰度。相同的三片段拼接产生 \(51\times90\) 的 fNIRS 张量，试验边界采用边缘填充。由于两个张量都包含 \(t+1\) 片段，该分支是离线估计器，而不是严格因果的实时解码器。

每种模态均由具有可学习邻接关系的多阶图网络编码，传播支持包括单位、一次和二次邻接关系。编码器生成 32 维通道嵌入。四头注意力沿 EEG 到 fNIRS 和 fNIRS 到 EEG 两个方向交换信息
\cite{vaswani2017attention,tsai2019multimodal}。对模态特异和跨模态 token 分别进行均值池化与最大池化，再输入多层回归器。训练目标为归一化标签的均方误差，加上 \(0.01\) 倍的对称对比对齐损失
\cite{chen2020simclr}。对齐项对同一样本的成对 EEG 与 fNIRS 表征进行正则化，但不主导回归目标。模型采用参与者分组折、AdamW 优化、较强 dropout，并依据折内验证 MSE 早停。

\subsection{偏重先验的固定融合}

最终估计为
\begin{equation}
 \hat y_{s,v,t}
 = w\odot \bar p_{v,t} + (1-w)\odot m_{s,v,t},
 \qquad w=[0.99,0.92].
\end{equation}
固定权重有意使视频—时间先验在效价和唤醒度上占主导。我们的目标是检验 EEG--fNIRS 能否为一个已经很强的低成本估计提供残差修正，而不是优化复杂的自适应门控模型。所有划分和两组评估都使用相同权重，使增量修正可以被直接解释。

\section{实验设置}
\label{sec:setup}

\subsection{数据集与参与者}

开发数据集包含 24 名参与者、所有参与者共同观看的 15 个熟悉视频、360 次试验和 36,864 个一赫兹效价—唤醒度目标
\cite{merpsdataset}。每次试验包含五秒静息基线。EEG 包含 64 个通道，fNIRS 包含 51 个通道上的六类信号。目标与预测均保留在原始 \([1,255]\) 尺度上。

\subsection{内部参与者留出评估}

参与者被划分为五折。对于每一折，视频—时间先验仅使用训练参与者重建，匹配的 EEG--fNIRS 检查点则预测留出参与者。全数据检查点不参与评估。所有折的预测与目标在计算样本汇总总体、效价和唤醒度 MAE 前进行拼接。该协议评估新观众对相同 15 个熟悉视频的反应。

\subsection{外部参与者独立评估}

外部队列包含 4 名参与者、60 次试验和 6,143 个一赫兹目标。这些参与者与开发集互不重叠。其目标不参与先验构建、检查点训练或推理，仅在本地评估时与预测连接。外部 EEG--fNIRS 估计对一个固定划分检查点和五个参与者分组折检查点取平均。固定融合将该集成结果与使用全部 24 名开发参与者构建的视频—时间先验相结合。

\subsection{来源变体}

表~\ref{tab:source-variants} 汇总了每种变体在推理时使用的信息。标签汇总列表示直接查找由开发参与者标签计算的缓存统计量；被评估参与者始终不进入这些汇总量。

\begin{table}[H]
  \caption{各来源变体在推理时使用的信息。标签汇总量由开发参与者标签计算；所有生理模型仍使用开发目标训练。}
  \label{tab:source-variants}
  \centering
  \scriptsize
  \setlength{\tabcolsep}{3.1pt}
  \begin{tabular}{@{}lcccc@{}}
    \toprule
    变体 &
    \shortstack{参与者\\生理信号} &
    \shortstack{视频\\键} &
    \shortstack{时间\\键} &
    \shortstack{标签\\汇总} \\
    \midrule
    全局常数 & 否 & 否 & 否 & 是 \\
    视频身份 & 否 & 是 & 否 & 是 \\
    视频—时间先验 & 否 & 是 & 是 & 是 \\
    EEG--fNIRS 分支 & 是 & 否 & 否 & 否 \\
    固定融合 & 是 & 是 & 是 & 是 \\
    \bottomrule
  \end{tabular}
\end{table}

全局常数是全部开发标签的二维中位数。视频身份变体使用每条开发视频轨迹的时间中位数，并在视频内保持不变。视频—时间先验使用完整中位数轨迹，并进行半径为三的视频内平滑。EEG--fNIRS 分支和固定融合分别遵循第~\ref{sec:method} 节及上述外部集成定义。内部档案包含最后三种变体，外部评估则包含全部五种变体。

\subsection{指标与报告方式}

我们在完全相同的预测键上报告样本汇总总体 MAE，并分别报告效价和唤醒度 MAE。总体 MSE 作为次要的外部诊断指标保留，而 MAE 始终是首要评价指标。对于外部异质性分析，成对误差还按参与者、熟悉视频及归一化视频内时间的十个等分区间进行汇总。由于外部队列较小，且相邻的一秒目标具有时间相关性，这些分组汇总仅作描述性分析。内部档案仅包含组件汇总指标，不含成对样本级残差。

\section{结果与分析}
\label{sec:results}

\subsection{RQ1：主要 MAE 表现}

表~\ref{tab:main-mae} 汇总了内部和外部结果。固定融合在两组评估中均取得最低总体 MAE，内部为 29.01，外部为 27.72。视频—时间先验与其分别仅相差 0.05 和 0.32 MAE，而 EEG--fNIRS 分支的误差明显更高。效价和唤醒度也呈现相同排序。

\begin{figure}[H]
  \centering
  \includegraphics[width=\columnwidth]{source_dominance.pdf}
  \caption{内部参与者留出 MAE。（a）EEG--fNIRS 分支、视频—时间先验和固定融合的总体 MAE。（b）按目标划分的从先验到融合的额外 MAE 降幅。}
  \Description{双面板图比较内部总体 MAE，以及从视频—时间先验到固定融合的较小额外降幅。}
  \label{fig:source-dominance}
\end{figure}

图~\ref{fig:source-dominance} 直观展示了相同的内部模式。从 EEG--fNIRS 分支到最终预测器的大幅误差降低几乎全部由视频—时间先验实现，而融合只提供了最后一个最小步骤。

\begin{table*}[t]
  \caption{原始标签尺度上的内部与外部 MAE。数值越低越好；短横线表示内部档案中没有对应来源变体。}
  \label{tab:main-mae}
  \centering
  \small
  \setlength{\tabcolsep}{6pt}
  \renewcommand{\arraystretch}{1.08}
  \begin{tabular}{@{}llrrrrr@{}}
    \toprule
    评估 & 目标 & 全局 & 视频身份 & 视频—时间先验 & EEG--fNIRS & 融合 \\
    \midrule
    内部 & 总体 & -- & -- & 29.06 & 47.35 & \textbf{29.01} \\
         & 效价 & -- & -- & 26.67 & 51.35 & \textbf{26.66} \\
         & 唤醒度 & -- & -- & 31.46 & 43.35 & \textbf{31.37} \\
    \midrule
    外部 & 总体 & 44.33 & 33.36 & 28.04 & 42.75 & \textbf{27.72} \\
         & 效价 & 47.46 & 32.56 & 25.34 & 45.29 & \textbf{25.20} \\
         & 唤醒度 & 41.20 & 34.17 & 30.74 & 40.21 & \textbf{30.25} \\
    \bottomrule
  \end{tabular}
\end{table*}

\subsection{RQ2：来源分解}

阶梯式外部变体确定了误差降低的来源。将全局常数替换为视频身份后，总体 MAE 降低 10.97 点；加入视频内时间后又降低 5.32 点。因此，共享反应同时包含视频特异性情感水平和显著的时间分辨轨迹。

在从 EEG--fNIRS 分支到固定融合的完整路径上，视频—时间先验解释了内部误差降低的 99.73\% 和外部误差降低的 97.89\%。固定融合仍然最优，但它相对于 EEG--fNIRS 分支的大部分优势已经可以从视频身份和视频内时间获得。图~\ref{fig:external-decomposition}a 展示了全部五种外部变体的完整排序。

\begin{figure*}[t]
  \centering
  \includegraphics[width=\textwidth]{external_source_decomposition.pdf}
  \caption{外部 MAE 与残差分析。（a）五种来源变体的总体 MAE。（b、c）参与者层面和视频层面的先验减融合增量；正值表示融合更优。（d）归一化视频时间区间上的 MAE。}
  \Description{四面板图展示外部总体误差、参与者层面和视频层面的成对融合增量，以及归一化视频时间上的误差。}
  \label{fig:external-decomposition}
\end{figure*}

\subsection{RQ3：参与者、视频与时间上的异质性}

外部残差表明，EEG--fNIRS 修正并不均匀。融合在 4 名参与者中的 3 名和 15 个熟悉视频中的 9 个上降低了 MAE（图~\ref{fig:external-decomposition}b,c）。参与者层面的变化范围为恶化 0.06 点至改善 0.69 点；视频层面的变化范围为恶化 0.30 点至改善 1.31 点。因此，更低的融合总体 MAE 与局部增益和损失同时存在。

视频内时间的贡献也会随视频进程变化（图~\ref{fig:external-decomposition}d）。在第一个归一化时间区间，视频身份达到 45.53 MAE，而视频—时间先验达到 13.40。该优势随时间推移而缩小，并且视频身份在最后五个区间中的每一个都略优。因此，时间带来的总体增益集中于特定阶段，而不是均匀分布。

\section{讨论与局限}
\label{sec:discussion}

\subsection{视频—时间先验为何有效}

视频—时间先验之所以有效，是因为部署场景保留了两个强坐标：熟悉视频身份和播放时间。参与者反应各异，但共享的视听序列仍会诱发群体层面轨迹。中位数汇总限制个体极端值，视频内平滑则在不跨越视频边界的前提下降低局部变化。因此，该预测器的优势来自数据结构，而不是模型复杂度。

这一特性赋予先验实际价值。既有观众建立群体轨迹后，无需采集 EEG 或 fNIRS，就能为新观众提供低成本的近似曲线。这类曲线可作为视频剪辑、广告投放、内容推荐或新观众预期反应初始化所需的群体反应基线。

\subsection{EEG--fNIRS 何时增加价值}

当参与者对群体轨迹的特异性偏差既可测量又稳定时，EEG--fNIRS 最有价值。外部残差表明，这一条件对部分参与者和视频成立，但并非全部。该异质性使传感成本成为设计决策的一部分：融合提供最低总体 MAE，而先验无需新的生理测量即可提供大部分收益。

限制修正能力的机制可能包括传感器噪声、一秒手工特征、不完善的血流动力学时序、跨参与者偏移，以及模型未显式分离相对先验的残差。当前分析定位了增益与损失，但无法区分这些解释。未来研究应检验更强的残差建模与显式血流动力学对齐。

\subsection{对实践与评估的启示}

结果也澄清了评估问题。保留熟悉视频的参与者留出评估衡量向新参与者的泛化，而视频留出评估衡量向新内容的迁移。未来研究应通过参与者—视频交叉划分分离这两个轴，并在每个外层训练划分内重建先验、特征标准化、检查点选择和融合权重
\cite{varma2006bias,varoquaux2017decoders}。参与者宏平均、视频宏平均和成对残差汇总还可以保留单一样本汇总分数所掩盖的层级结构。

在部署时，应根据残差改善相对于采集成本的价值，在先验与融合之间进行选择。先验为熟悉内容提供简单的群体基线，而 EEG--fNIRS 可保留给值得通过额外设备进行参与者特异性修正的场景。

\subsection{局限与未来工作}

本文结论受熟悉视频场景限制：所有参与者观看相同的 15 个视频，因此没有测量向未见内容的迁移。外部队列仅包含 4 名参与者，且相邻的一秒目标具有相关性；因此，异质性分析仅是描述性的，不能视为稳定的参与者层面推断。融合权重与检查点选择未通过完全嵌套流程确定；内部和外部融合增量也不能直接比较，因为二者使用不同的生理集成和留出群体。EEG--fNIRS 分支使用一秒未来上下文，且未显式建模 fNIRS 血流动力学时移，因此是离线估计器。本文评估预测误差，而非下游应用效用。匿名标识阻碍人口统计学或亚组分析，当前生理架构也仅代表一个比较器。参与者—视频交叉评估、嵌套选择和更强的残差建模，是扩大验证范围的主要方向。

\newpage
\balance
\section{结论}
\label{sec:conclusion}

本文构建了排除被评估参与者标签的折内视频—时间先验，作为熟悉视频连续情感回归的强低成本基线。与 EEG--fNIRS 分支固定融合后，模型在内部和外部评估中均取得最低 MAE；但来源分解表明，大部分误差降低来自视频身份和视频内时间。EEG--fNIRS 仅提供较小且异质的残差修正。该方法可在熟悉内容可用时支持群体反应曲线；更广泛的结论仍需参与者—视频交叉评估和嵌套残差模型选择。
